\documentclass[12pt,a4paper]{article}
\usepackage[utf8]{inputenc}
\usepackage[brazilian]{babel}
\usepackage[left=3cm,top=3cm,right=2cm,bottom=2cm]{geometry}
\usepackage{graphicx}
\usepackage{setspace}
\usepackage[colorlinks=true, urlcolor=blue, linkcolor=blue]{hyperref}
\usepackage{xurl}
\usepackage{tikz}
\usepackage{listings}
\usepackage{booktabs}
\usepackage{float}

\onehalfspacing

\begin{document}

\begin{titlepage}    
    \begin{center}     
        {\large \textbf{UNIVERSIDADE DE SÃO PAULO}} \\
        {\large \textbf{INSTITUTO DE MATEMÁTICA, ESTATÍSTICA E CIÊNCIA DA COMPUTAÇÃO}} \\

        \vspace{5cm}
        \vfill
        
        \begin{minipage}{0.8\textwidth}
            \centering
            {\Large \textbf{MAC0216 - Trabalho Prático 3 -  Decodificador UTF-8 em Assembly}} \\
            \vspace{0.8cm}
            {\large Gabriela Kwong Yamaguti (16967502),\\Isabela Suemi Aoki (15637975),\\Yoonsik Hur (16905558),\\Viviane Aiko Yamamoto (5515726)}
        \end{minipage}
        
        \vfill
        \vfill
        
        {\large São Paulo} \\
        {\large 2026}
    \end{center}
    
    \vspace*{\fill}
\end{titlepage}

\newpage
\tableofcontents
\thispagestyle{empty}

\newpage
\pagenumbering{arabic}

% 1. Introdução -----------------------------------------------------
\section{Introdução}

O presente trabalho teve como objetivo desenvolver um programa em Assembly NASM para Linux de 64 bits. Ele lê um arquivo de texto em UTF-8 e produz outro arquivo com o código Unicode de cada caractere, escrito em hexadecimal. O desenvolvimento foi versionado com Git em um repositório no GitHub, e o programa pode ser montado e ligado com o Makefile incluído.

\subsection{Explicação do Funcionamento do Programa}
\subsubsection*{Convenção de chamada}
Todos os parâmetros das funções são passados pela pilha, ocupando 8 bytes cada, empilhados do último para o primeiro. Dentro da função, o primeiro parâmetro fica em \texttt{[rbp+16]} e o segundo em \texttt{[rbp+24]}. O chamador remove os parâmetros com \texttt{add rsp}, e o valor de retorno é entregue em \texttt{RAX}.

\subsubsection*{Funções obrigatórias}
\begin{enumerate}
    \item \texttt{escreve\_string}: conta os caracteres até o NUL e escreve na saída padrão (\texttt{sys\_write}).
    \item \texttt{le\_string}: lê da entrada padrão (\texttt{sys\_read}) byte a byte até o ENTER ou até \texttt{tam\_max - 1}. Grava o NUL e devolve em \texttt{RAX} a quantidade lida.
    \item \texttt{abre\_arquivo}: valida o modo (0 ou 577) e chama \texttt{sys\_open} com permissão 438. Devolve o descritor ou $-1$ em caso de erro.
    \item \texttt{fecha\_arquivo}: fecha o descritor com \texttt{sys\_close}. Devolve 0 ou $-1$.
    \item \texttt{le\_byte\_arquivo}: lê 1 byte do arquivo para o endereço dado. Devolve em \texttt{RAX} a quantidade lida (0 indica fim do arquivo).
    \item \texttt{grava\_string\_arquivo}: calcula o tamanho da string até o NUL e grava no arquivo, incluindo o LF e excluindo o NUL.
    \item \texttt{gera\_string\_hexadecimal}: converte o valor para hexadecimal maiúsculo com prefixo \texttt{0x}, seguido de LF e NUL. 
\end{enumerate}

\subsubsection*{Fluxo do programa principal}
\begin{enumerate}
    \item Solicita e lê os nomes dos arquivos de entrada e de saída, com \texttt{escreve\_string} e \texttt{le\_string}.
    \item Abre o arquivo de entrada para leitura (modo 0) e o de saída para escrita, criação e truncamento (modo 577), com \texttt{abre\_arquivo}. Se qualquer abertura retornar $-1$, o programa encerra com código 1.
    
    \item Entra no laço principal, que lê um byte por vez com \texttt{le\_byte\_arquivo} até o retorno ser 0 (que representa o fim do arquivo). Para cada caractere:
    \begin{enumerate}
        \item \textbf{Identifica o tamanho da sequência}: com \texttt{AND} e \texttt{CMP} sobre o primeiro byte. O registrador \texttt{R13} guarda o número de bytes de continuação e \texttt{AL} armazena apenas os bits de dados do primeiro byte. 
        \item \textbf{Reconstrói o code point} em \texttt{R12}: para cada byte de continuação, lê o byte, mantém seus 6 bits de dados (\texttt{AND 00111111b}), desloca \texttt{R12} 6 bits à esquerda (\texttt{SHL}) e junta os bits com \texttt{OR}.
        \item \textbf{Converte para hexadecimal:} usando \texttt{gera\_string\_hexadecimal}, o valor é dividido por 16 repetidamente e os restos são empilhados. Ao desempilhar, a ordem fica invertida e cada resto é convertido em dígito (\texttt{0}--\texttt{9}: soma 48; \texttt{A}--\texttt{F}: soma 55). O buffer recebe \texttt{0x}, os dígitos, e depois, LF e NUL.
        \item \textbf{Grava a linha} no arquivo de saída com \texttt{grava\_string\_arquivo}.
    \end{enumerate}
    \item Ao chegar ao fim do arquivo, fecha os dois descritores com \texttt{fecha\_arquivo} e encerra com \texttt{sys\_exit} (código 0).
\end{enumerate}

% 2. Testes Realizados -----------------------------------------------
\section{Testes Realizados}

% 3. Dificuldades Encontradas -----------------------------------------
\section{Dificuldades Encontradas}

% 4. Contribuição dos Integrantes -----------------------------------------
\section{Contribuição dos Integrantes}
\begin{enumerate}
    \item Gabriela:
    \item Isabela: desenvolvimento de \texttt{gera\_string\_hexadecimal} e do fluxo principal do programa, além da escrita do relatório.
    \item Viviane:
    \item Yoonsik:
\end{enumerate}

% 5. Links -----------------------------------------
\section{Links}
Link do GitHub: \url{https://github.com/isaaoki/mac0216-decodificador-assembly}

\end{document}